\documentclass[JEP,XML,SOM,Unicode,NoEqCountersInSection,published]{cedram}
\datereceived{2024-05-25}
\dateaccepted{2025-06-26}
\dateepreuves{2025-07-01}

\usepackage[scr=rsfs,cal=euler]{mathalfa}
\multlinegap0pt

\hyphenation{dimen-sion remark remarks}

\let\ra\rightarrow
\DeclareMathOperator{\supp}{supp}


\newcommand{\RR}{\mathbb R}
\newcommand{\NN}{{\mathbb N}}
\renewcommand{\Im}{\mathop{\mathrm{Im}}}

\newtheorem{theorem}{Theorem}
\newtheorem{lemma}{Lemma}
\newtheorem{proposition}{Proposition}
\newtheorem{corollary}{Corollary}

\theoremstyle{definition}
\newtheorem{remark}{Remark}
\newtheorem{definition}{Definition}
\newtheorem{example}{Example}

\datepublished{2025-07-02}
\begin{document}
\frontmatter
\title{A holographic global uniqueness in passive\nobreakspace imaging}

\author[\initial{R.} \middlename{G.} \lastname{Novikov}]{\firstname{Roman} \middlename{G.} \lastname{Novikov}}
\address{CMAP, CNRS, École polytechnique,\\
Institut Polytechnique de Paris, Palaiseau, France}
\address{\& IEPT RAS, 117997 Moscow, Russia}
\email{novikov@cmap.polytechnique.fr}
\urladdr{http://www.cmapx.polytechnique.fr/~novikov}

\begin{abstract}
We consider a radiation solution $\psi$ for the Helmholtz equation in an exterior region in $\mathbb{R}^3$. We show that the restriction of $\psi$ to any ray $L$ in the exterior region is uniquely determined by its imaginary part $\Im\psi$ on an interval of this ray. As a corollary, the restriction of $\psi$ to any plane $X$ in the exterior region is uniquely determined by $\Im\psi$ on an open domain in this plane. These results have holographic prototypes in the recent work Novikov (2024, Proc.\ Steklov Inst.\ Math.\ 325, 218-223). In particular, these and known results imply a holographic type global uniqueness in passive imaging and for the Gelfand-Krein-Levitan inverse problem (from boundary values of the spectral measure in the whole space) in the monochromatic case. Some other surfaces for measurements instead of the planes $X$ are also considered.
\end{abstract}

\subjclass{35J05, 35J08, 35P25, 35R30}

\keywords{Helmholtz equation, Schrödinger equation, radiation solutions, Gelfand-Krein-Levitan problem, passive imaging,
holographic global uniqueness}

\altkeywords{Équation de Helmholtz, équation de Schrödinger, solutions de rayonnement, problème de Gelfand-Krein-Levitan, imagerie passive, unicité globale holographique}

\alttitle{Une unicité globale holographique en imagerie passive}

\begin{altabstract}
Nous considérons une solution de rayonnement $\psi$ pour l'équation de Helmholtz dans une région extérieure de $\mathbb R^3$. Nous montrons que la restriction de $\psi$ à tout rayon $L$ de la région extérieure est déterminée de manière unique par sa partie imaginaire $\Im\psi$ sur un intervalle de ce rayon. En corollaire, la restriction de $\psi$ à tout plan $X$ de la région extérieure est déterminée de manière unique par $\mathop{\mathrm{Im}}\psi$ sur un domaine ouvert de ce plan. Ces résultats ont des prototypes holographiques dans l'article récent de Novikov (2024, Proc.\ Steklov Inst.\ Math.\ 325, 218-223). En particulier, ces résultats et des résultats connus impliquent une unicité globale de type holographique en imagerie passive et pour le problème inverse de Gelfand-Krein-Levitan (à partir des valeurs au bord de la mesure spectrale dans l'espace entier) dans le cas monochromatique. D'autres surfaces de mesure que les plans $X$ sont également considérées.
\end{altabstract}

\maketitle
\tableofcontents
\mainmatter

\section{Introduction}\label{sec:1}

We consider the Helmholtz equation
\begin{equation}
\label{eq:1.1}
-\Delta\psi(x) = \kappa^2 \psi(x),\quad x\in \mathcal{U},\ \kappa>0,
\end{equation}
where $\Delta$ is the Laplacian in $x$, and $\mathcal U$ is an exterior region in $\RR^3$
that is $\mathcal U$ is an open unbounded connected set in $\RR^3$ consisting of all points in exterior of a closed bounded regular surface $S$ (as in \cite{A}, \cite{W}).

For equation \eqref{eq:1.1} we consider the radiation solutions $\psi$ that is the solutions satisfying the Sommerfeld's radiation condition
\begin{equation}\label{eq:1.2}
|x|\Bigl(\frac{\partial}{\partial|x|}-i\kappa\Bigr)\psi(x)\to 0\quad \mbox{as} \ |x|\ra +\infty,
\end{equation}
uniformly in $x/|x|$. We assume that $\psi \in C^2$ in the closure of $\mathcal{U}$.

Let
\begin{equation}\label{eq:1.3}
L=L_{x_0, \theta}=\{x\in\RR^3:\ x=x(s)=x_0+s\theta,\ 0<s <+\infty \}, \quad x_0\in\RR^3, \ \theta\in\mathbb S^2,
\end{equation}
where $\mathbb S^2$ is the unit sphere in $\RR^3$.

In the present work we show that, for any radiation solution $\psi$ and any ray $L\subset ~\mathcal{U}$,
the restriction of $\psi$ to $L$ is uniquely determined by its imaginary part $\Im\psi$ on an arbitrary interval of $L$;
see Theorem \ref{th:1} in Section \ref{sec:2}.

As a corollary, we also obtain that, for any radiation solution $\psi$ and any plane $X\subset \mathcal{U}$, the restriction of $\psi$ to $X$ is uniquely determined by $\Im\psi$
on an arbitrary open domain of $X$; see Theorem \ref{th:2} in Section \ref{sec:2}.

As a further corollary, we obtain that, for any radiation solution $\psi$ and any plane $X\subset \mathcal{U}$, the solution $\psi$ on the whole $\mathcal{U}$ is also uniquely determined by $\Im\psi$
on an arbitrary open domain of $X$; see Corollary \ref{cor:1} in Section \ref{sec:2}.

These results have holographic prototypes in the recent work \cite{N5}, where the aforementioned reconstructions of $\psi$ are considered from the intensity $|e^{ikx}+\psi|^2$ in place of $\Im\psi$.
Here, $e^{ikx}$ is a plane wave solution of \eqref{eq:1.1}, \ie $k\in\RR^3,\ |k|=\kappa$. The results of~\cite{N5} solve one of old mathematical questions of holography and
admit straightforward applications to phaseless inverse scattering.

In the present work our studies are motivated by the Gelfand-Krein-Levitan inverse problem (from boundary values of the spectral measure in the whole space)
and by passive imaging; see, e.g., \cite{AHN1}, \cite{AHN2}, \cite{Ber}, \cite{BSS}, \cite{GP}, \cite{G}, \cite{MH}, \cite{WL}.

The Gelfand-Krein-Levitan problem in question (in its fixed energy version in dimension $d=3$) consists in determining the potential $v$ in the Schrödinger equation
\begin{equation}\label{eq:1.4}
-\Delta\psi(x) +v(x)\psi(x)= \kappa^2 \psi(x) +\delta(x-y),\ x,y\in \RR^3,\ \kappa>0,
\end{equation}
from the imaginary part of the outgoing Green function $R_v^+(x, y, \kappa)$ for one $\kappa$ and all $x$, $y$ on some part of the boundary of a domain containing the support of $v$.
Here, $\delta$ is the Dirac delta function.
In this problem $\Im R_v^+$ is related to the spectral measure of the Schrödinger operator $H=-\Delta+v$. More precisely, $H$
admits the following spectral decomposition in $L^2(\RR^d)$, at least, for real-valued compactly supported $v\in L^{\infty}(\RR^d)$:
\begin{equation}\label{eq:1.6}
H=\int_0^{\infty}\kappa^2 d\mu_{\kappa}+\sum_{j=1}^N E_j \pi_j, \quad d\mu_{\kappa}=\frac{2}{\pi}\Im R_v^+(\kappa)\kappa d\kappa,
\end{equation}
where $d\mu_{\kappa}$ is the positive part of the spectral measure for $H$, $E_j$ are nonpositive eigenvalues for $H$ and
$\pi_j$ are orthogonal projectors on corresponding eigenspaces, $R_v^+(\kappa)=(H-\kappa^2-i0)^{-1}$ is the limiting absorption resolvent for $H$,
whose Schwartz kernel is given by $R_v^+(x, y, \kappa)$; see, e.g., \cite[Lem.\,14.6.1]{H}.

Note that the terminology ``Gelfand-Krein-Levitan problem'' is not conventional.
Our motivation for this terminology is based on the Yu.\,M.\,Berezanskii's work \cite{Ber},
which is one of the very first mathematical works on the multidimensional inverse problems for differential equations.
According to \cite{Ber}, the problem of finding potential $v$ in the multidimensional Schrödinger equation
from boundary values of the spectral measure for some fixed boundary condition was formulated originally
by M.\,G.\,Krein, I.\,M.\,Gelfand and B.\,M.\,Levitan at a conference on differential equations in Moscow in 1952.
In addition, \cite{Ber} gives, in particular, uniqueness theorems on such problems including the case of the problem of determining $v$
from boundary values (together with some normal derivatives) of the spectral measure arising in \eqref{eq:1.6} for all real energies on a part of the boundary,
at least, for piecewise real-analytic $v$.

In addition, $R_v^+(x, y, \kappa)$, for fixed $y$, can be defined as the solution $\psi$ of equation~\eqref{eq:1.4} with the radiation condition \eqref{eq:1.2}.
For more details, see \cite{AHN1}, \cite{Ber} and references therein.

Equation \eqref{eq:1.4} at fixed $\kappa$ can be also considered as the Helmholtz equation of acoustics or electrodynamics for monochromatic waves,
where complex-valued $v(x)= v(x,\kappa)$ is related to the perturbation of the refraction index.
In particular, the aforementioned mathematical problem of recovering $v$ from boundary values of $\Im R_v^+$ arises in the framework of
passive acoustic tomography (in ultrasonics, ocean acoustics, local helioseismology). In these framework
$\Im R_v^+$ is related to cross correlations of wave fields generated by random sources; see, e.g., formula (49) in \cite{G}.
For more details, see \cite{AHN1}, \cite{AHN2}, \cite{BSS}, \cite{G}, \cite{S}, \cite{WL} and references therein.

Note that
\begin{equation}\label{eq:1.7}
R_0^+(x, y, \kappa)=\frac{e^{i\kappa|x-y|}}{4\pi|x-y|},\quad x,y\in \RR^3,
\end{equation}
where $R_0^+$ is the outgoing Green function for equation \eqref{eq:1.4} with $v \equiv 0$.

Let
\begin{equation}\label{eq:1.8}
R_{v,sc}^+(x, y, \kappa)=R_v^+(x, y, \kappa)-R_0^+(x, y, \kappa),\quad x,y\in \RR^3.
\end{equation}

Suppose that
\begin{equation}
\label{eq:1.5}
\supp v \subset \RR^3\setminus \overline {\mathcal{U}},
\end{equation}
where $\overline {\mathcal{U}}$ is the closure of $\mathcal{U}$.
Then, in view of the definitions of $R_{v,sc}^+$, $R_v^+$, and $R_0^+$, the function $\psi=R_{v,sc}^+(x, y, \kappa)$ is a radiation solution of equation \eqref{eq:1.1}
for each $y\in \RR^3$.

Therefore, the aforementioned results on recovering a radiation solution $\psi$ from $\Im \psi$ give a reduction of the Gelfand-Krein-Levitan problem
(of inverse spectral theory and passive imaging in dimension $d=3$) to the inverse scattering problem of finding~$v$ in \eqref{eq:1.4} from boundary values of $R_v^+$.
Note that studies on the latter problem also go back to \cite{Ber}.

This reduction and known results imply, in particular, that $v$ in \eqref{eq:1.4} is uniquely determined by $\Im R_v^+(x, y, \kappa)$
for one $\kappa$ and all $x$, $y$ on an arbitrary open domain $\mathcal{D}$ of $X$,
where $\supp v \subset \Omega$, $X$ is a plane in $\RR^3 \setminus \overline {\Omega}$, $\Omega$ is an open bounded connected domain in $\RR^3$, $\overline {\Omega}$ is the closure of $\Omega$;
see Theorem~\ref{th:3} in Section \ref{sec:2}.
By this result we continue studies of \cite{Ber} mentioned above and relatively recent studies of \cite{AHN1} and \cite{AHN2}.

We also consider other surfaces for measurements
instead of the planes~$X$; see Example \ref{ex:1} and Theorems \ref{th:4} and \ref{th:5} in Section \ref{sec:2}.

The main results of this work are presented in more detail and proved in Sections~\hbox{\ref{sec:2}--\ref{sec:7}}.

\subsubsection*{Acknowledgements}
The author thanks the referees for remarks that have helped to improve the presentation.

\section{Main results}\label{sec:2}

In this work our key result is as follows.

\begin{theorem}\label{th:1}
Let $\psi$ be a radiation solution of equation \eqref{eq:1.1} as in \eqref{eq:1.2}.
Let $L$ be a ray as in \eqref{eq:1.3} such that $L \subset \mathcal{U}$, where $\mathcal{U}$ is the region in \eqref{eq:1.1}. Then $\psi$ on $L$ is uniquely determined by $\Im \psi$ on $\Lambda$,
where $\Lambda$ is an arbitrary non-empty open interval of~$L$.
\end{theorem}

As a corollary, we also get, in particular, the following result.

\begin{theorem}\label{th:2}
Let $\psi$ be a radiation solution of equation \eqref{eq:1.1} as in \eqref{eq:1.2}.
Let $X$ be a two-dimensional plane in $\RR^3$ such that $X\subset \mathcal{U}$.
Then $\psi$ on $X$ is uniquely determined by $\Im \psi$ on $\mathcal{D}$,
where $\mathcal{D}$ is an arbitrary non-empty open domain of $X$.
\end{theorem}

Theorem \ref{th:1} is proved in Section \ref{sec:4} using the Atkinson-Wilcox expansion of \cite{A}, \cite{W} for the radiation solutions $\psi$ of equation \eqref{eq:1.1}, and a modified version of holographic techniques of \cite{N1}, \cite{N5}.
In particular, in this proof we use Proposition~\ref{prop:1} of Section~\ref{sec:3}, which yields a two-point approximation for $\psi$
in terms of $\Im \psi$. This two-point approximation is also of independent interest.

Note that $\psi$ and $\Im\psi$ are real-analytic on $\mathcal{U}$, and, therefore, on $L$ in Theorem \ref{th:1} and on $X$ in Theorem \ref{th:2}.
Because of this analyticity, Theorem \ref{th:1} reduces to the case when $\Lambda=L$ and Theorem \ref{th:2} reduces to the case when $\mathcal{D}=X$.

Using this reduction, Theorem \ref{th:2} is proved as follows.
We assume that $\mathcal{D}=X$.
Then to determine $\psi$ at an arbitrary $x\in X$, we consider a ray $L=L_{x_0, \theta}\subset X$ such that $x\in L$ and use Theorem \ref{th:1} for this $L$.

\begin{corollary}\label{cor:1}
Under the assumptions of Theorem \ref{th:2}, the imaginary part $\Im\psi$ on~$\mathcal{D}$, uniquely determines $\psi$
in the entire region $\mathcal{U}$.
\end{corollary}

Corollary \ref{cor:1} follows from Theorem \ref{th:2}, formula \eqref{eq:3.8a} recalled in Section~\ref{subsec:3.3}, and analyticity of $\psi$ in $\mathcal{U}$.

\begin{remark}\label{rem:1}
In the one dimensional case, an analog of Theorem \ref{th:1} follows from a very simple form of one dimensional radiation solutions.
In particular, in this case an analog of the two-point approximation of Proposition \ref{prop:1} is exact.
The two dimensional case is considered in \cite{NN2} using results of \cite{K}, \cite{NN1} in place of the three dimensional Atkinson-Wilcox expansion \eqref{eq:3.2} used in the present work.
We expect that the case of dimension $d>3$ is similar to the three-dimensional case if $d$ is odd and is similar to the two-dimensional case if $d$ is even.
\end{remark}

Theorems \ref{th:1} and \ref{th:2}, and Corollary \ref{cor:1} have holographic prototypes in \cite{N5}; see Introduction for some comments in this connection.

Using Theorem \ref{th:2} and known results on direct and inverse scattering we obtain the following global uniqueness theorem for the Gelfand-Krein-Levitan inverse problem mentioned in Introduction.

\begin{theorem}\label{th:3}
Let $v \in L^{\infty}(\RR^3)$, $\supp v \subset \Omega$, and $X \subset \RR^3 \setminus \overline {\Omega}$, where $\Omega$ is an open bounded connected domain in $\RR^3$, $\overline {\Omega}$ is the closure of $\Omega$,
and $X$ is a two-dimensional plane. Let property \eqref{eq:3.11} hold for fixed $\kappa>0$ and $R_v^+(x, y, \kappa)$ be the outgoing Green function for equation \eqref{eq:1.4}.
Then $v$ is uniquely determined by $\Im R_v^+(\cdot, \cdot, \kappa)$ on $\mathcal{D} \times \mathcal{D}$,
where $\mathcal{D}$ is an arbitrary non-empty open domain of $X$.
\end{theorem}
In Theorem \ref{th:3} we do not assume that $v$ is real-valued,
but we assume that property~\eqref{eq:3.11} formulated in Section \ref{subsec:3.4} holds.

\begin{remark}\label{rem:2}
Under the conditions that $v$ is real-valued, $v \in L^{\infty}(\RR^d)$, $\supp v \subseteq \overline {\Omega}$,
where $\Omega$ is an open bounded domain in $\RR^d$ with $\partial\Omega \in C^{2,1}$, $\overline {\Omega}=\Omega \cup \partial\Omega$, $d \geq 2$,
a~local uniqueness for determining $v$ from
$\Im R_v^+(\cdot, \cdot, \kappa)$ on $\partial\Omega \times \partial\Omega$ for one $\kappa$
is proved in~\cite{AHN1}. The assumption that $v$ is real-valued is very essential in this proof.
\end{remark}

\begin{remark}\label{rem:3}
In connection with applications to helioseismology, it is natural to assume that $v$ is complex-valued, $v \in L^{\infty}(\RR^3)$, $v(x)=\tilde v(|x|)$, $\tilde v(r)=\alpha/r$, $r \geq r_0$,
for some constants $\alpha \in \RR$ and $r_0>0$. Let $M_r=\mathbb S^2_r \times \mathbb S^2_r$ and $r_2>r_1>r_0$,
where $\mathbb S^2_r$ is defined by \eqref{eq:2.1}.
Under these assumptions, a global uniqueness for determining~$v$ from
$\Im R_v^+(\cdot, \cdot, \kappa)$ on $M_{r_1} \cup M_{r_2}$, for fixed $\kappa$ and nonsingular pair $r_1$, $r_2$,
is proved in~\cite{AHN2}.
\end{remark}

One can see that
Theorem \ref{th:3} contains a principal progress on the Gelfand-Krein-Levitan inverse problem in comparison with the results of \cite{AHN1}, \cite{AHN2} mentioned in Remarks \ref{rem:2} and \ref{rem:3}.
In comparison with the aforementioned result of \cite{AHN1}, Theorem \ref{th:3} is global and does not assume that $v$ is real-valued.
In comparison with the aforementioned result \cite{AHN2}, Theorem \ref{th:3} does not assume that $v$ is spherically symmetric.

Theorem \ref{th:3} is proved in Section \ref{sec:5}.
In particular, in this proof we use Proposition~\ref{prop:2} of Section \ref{sec:5},
which stays that $v$ is uniquely determined by $R_v^+(\cdot, \cdot, \kappa)$ on $X\times X$ for fixed~$\kappa$.
To our knowledge, Proposition \ref{prop:2} is slightly different from the results reported in the literature.
Therefore, just in case, for completeness of presentation this proposition is also proved in Section \ref{sec:5}.

In connection with other surfaces of measurements instead of the planes $X$ our results are as follows.

The results of Theorems \ref{th:1} and \ref{th:2} don't hold for some other curves in place of $L$ and surfaces in place of $X$.
An example is as follows.

Let
\begin{equation}
\label{eq:2.1}
\mathbb S^2_r=\{x\in\RR^3:\ |x|=r\},\quad r>0.
\end{equation}

\begin{example}\label{ex:1}
Let $\psi(x)= G^+(x,\kappa)$, where $G^+$ is defined by \eqref{eq:3.8b}. Then $\psi$ is a non-zero radiation solution of equation \eqref{eq:1.1} if $0 \in \RR^3\setminus \overline {\mathcal{U}}$, but $\Im\psi \equiv 0$ on the spheres~$\mathbb S^2_r$ for
$r=n\pi/ \kappa$, $n \in\NN$.
\end{example}
Nevertheless, the uniqueness results of Theorem \ref{th:2}, Corollary \ref{cor:1}, and Theorem \ref{th:3} remain valid for many interesting surfaces instead of the planes $X$.
Suppose that
\begin{equation}\label{eq:2.2}
\begin{aligned}
&\mathcal{B}\ \mbox{is\ an\ open\ bounded\ domain\ in}\ \RR^3, \\
&Y=\partial \mathcal{B}\ \mbox{is\ real-analytic\ and\ connected}.
\end{aligned}
\end{equation}

In particular, we have the following uniqueness theorems.

\begin{theorem}\label{th:4}
Let $\psi$ be a radiation solution of equation \eqref{eq:1.1} as in \eqref{eq:1.2}.
Let $Y$ be a surface as in \eqref{eq:2.2},
where $\kappa$ is not a Dirichlet eigenvalue for $\mathcal{B}$, and $\overline{\mathcal{B}}= \mathcal{B} \cup Y \subset \mathcal{U}$.
Then $\psi$ on $\mathcal{U}$ is uniquely determined by $\Im \psi$ on $\mathcal{D}$,
where $\mathcal{D}$ is an arbitrary non-empty open domain of $Y$.
\end{theorem}

\begin{theorem}\label{th:5}
Let $v \in L^{\infty}(\RR^3)$, $\supp v \subset \RR^3 \setminus \overline {\mathcal{U}}$, where $\mathcal{U}$ is as in \eqref{eq:1.1},
and property \eqref{eq:3.11} holds for fixed $\kappa>0$.
Let $Y$ be a surface as in \eqref{eq:2.2},
where $\kappa$ is not a Dirichlet eigenvalue for $\mathcal{B}$, and $\overline{\mathcal{B}}= \mathcal{B} \cup Y \subset \mathcal{U}$.
Then $v$ is uniquely determined by $\Im R_v^+(\cdot, \cdot, \kappa)$ on $\mathcal{D} \times \mathcal{D}$,
where $\mathcal{D}$ is an arbitrary non-empty open domain of $Y$.
\end{theorem}

Theorems \ref{th:4} and \ref{th:5} are proved in Sections \ref{sec:6} and \ref{sec:7}.
Note that the determinations in Theorems \ref{th:1}, \ref{th:2}, \ref{th:3} and especially in Theorems \ref{th:4} and \ref{th:5} include analytic continuations.
Studies on more stable reconstructions appropriate for numerical implementation will be continued elsewhere.
In this respect one can use, in particular, results of Section~\ref{subsec:3.2}.

\section{Preliminaries}\label{sec:3}

\subsection{The Atkinson-Wilcox expansion}

Let
\begin{equation}
\label{eq:3.1}
B_r=\{x\in\RR^3:\ |x|<r\},\quad r>0.
\end{equation}
If $\psi$ is a radiation solution of equation \eqref{eq:1.1} and $\RR^3\setminus B_r \subset \mathcal{U}$, then
the following Atkinson-Wilcox expansion holds:
\begin{equation}
\label{eq:3.2}
\psi(x)=\frac{e^{i\kappa|x|}}{|x|}\sum\limits_{j=1}^{\infty}\frac{f_j(\theta)}{|x|^{j-1}} \quad \text{for}\ x\in \RR^3\setminus B_r,\ \theta=\frac{x}{|x|},
\end{equation}
where the series converges absolutely and uniformly; see \cite{A}, \cite{W}.

\subsection{A two-point approximation for $\psi$}\label{subsec:3.2}

Let
\begin{equation}
\label{eq:3.3}
I(x)=|x|\Im\psi(x),\quad x\in\mathcal{U},
\end{equation}
where $\psi$ is a radiation solution of equation \eqref{eq:1.1}.

We have that
\begin{equation}
\label{eq:3.4}
2iI(x)=e^{i\kappa|x|}f_1(\theta)-e^{-i\kappa|x|} \overline {f_1(\theta)}+O(|x|^{-1})\quad \text{as} \ |x|\ra +\infty,
\end{equation}
uniformly in $\theta=x/|x|$, where $f_1$ is the leading coefficient in \eqref{eq:3.2}.

\begin{proposition}\label{prop:1}
Let $\psi$ be a radiation solution of equation \eqref{eq:1.1}. Then
\begin{equation}\label{eq:3.5}
\begin{aligned}
&f_1(\theta)= \frac{1}{\sin(\kappa\tau)}\bigl(-e^{-i\kappa|y|}I(x)+e^{-i\kappa|x|}I(y)+O(|x|^{-1})\bigr),\\
&x,y\in L_{x_0, \theta},\quad x_0=0,\ y=x+\tau\theta, \quad\theta\in\mathbb S^2,\quad \tau>0,
\end{aligned}
\end{equation}
uniformly in $\theta$, where $f_1$ is the coefficient in \eqref{eq:3.4}, $I$ is defined by \eqref{eq:3.3}, $L_{x_0, \theta}$ is the ray defined by \eqref{eq:1.3},
and $\sin(\kappa\tau)\ne 0$ for fixed $\tau$.
\end{proposition}

Formula \eqref{eq:3.5} is a two-point approximation for $f_1$ and together with \eqref{eq:3.2} also gives a two-point approximation for $\psi$
in terms of $I$.
For phaseless inverse scattering and holography formulas of such a type go back to \cite{N1} (see also \cite{NS2}, \cite{NSh}, \cite{N5}).

\begin{remark}\label{rem:4}
If an arbitrary function $I$ on $L_{0, \theta}$ satisfies \eqref{eq:3.4},
then formula \eqref{eq:3.5} holds, for fixed $\theta\in\mathbb S^2$.
\end{remark}

We obtain \eqref{eq:3.5} from the system of equations for $f_1$ and $\overline {f_1}$:
\begin{equation}\label{eq:3.6}
\begin{aligned}
&e^{i\kappa|x|}f_1(\theta)-e^{-i\kappa|x|} \overline {f_1(\theta)}=2iI(x)+O(|x|^{-1}),\\
&e^{i\kappa|y|}f_1(\theta)-e^{-i\kappa|y|} \overline {f_1(\theta)}=2iI(y)+O(|y|^{-1}),
\end{aligned}
\end{equation}
where $x$, $y$ are as in \eqref{eq:3.5}. In particular, we use that $|y|=|x|+\tau$ and that
\begin{equation}
\label{eq:3.7}
D=2i \sin(\kappa\tau),
\end{equation}
where $D$ is the determinant of system \eqref{eq:3.6}.

In turn, \eqref{eq:3.6} follows from \eqref{eq:3.4}.

\subsection{A Green type formula}\label{subsec:3.3}

The following formula holds:
\begin{align}
\psi(x)&= 2\int_{X}\frac{\partial G^+(x-y,\kappa)}{\partial \nu_{y}}\psi(y)dy,\quad x\in V_X, \label{eq:3.8a}\\
G^+(x,\kappa)&= -\frac{e^{i\kappa|x|}}{4\pi|x|},\quad x\in \RR^3, \label{eq:3.8b}
\end{align}
where $\psi$ is a radiation solution of equation \eqref{eq:1.1}, $X$ and $V_X$ are plane and open half-space in $\mathcal{U}$,
where $X$ is the boundary of $V_X$, $\nu$ is the outward normal to $X$ relative to~$V_X$; see, for example, \cite[Eq.\,(5.84)]{BR}.

Recall that $R_0^+(x, y, \kappa)=-G^+(x-y,\kappa)$, where $R_0^+$ is the outgoing Green function for equation \eqref{eq:1.4} with $v \equiv 0$.

For completeness of presentation note that formula \eqref{eq:3.8a} follows from the formula
\begin{equation}\label{eq:3.9a}
\begin{aligned}
\psi(x)&= \int_{X}\frac{\partial G^+_X(x,y,\kappa)}{\partial \nu_{y}}\psi(y)dy,\quad x\in V_X,\\
G^+_X(x,y,\kappa)&= G^+(x-y,\kappa)-G^+(x-y^*,\kappa),\quad x,y\in V_X \cup X,
\end{aligned}
\end{equation}
where $y^*$ is symmetric to $y$ with respect to $X$.
The point is that $G^+_X$ is the Green function for the Helmholtz operator $\Delta+\kappa^2$ in $V_X$
with Dirichlet boundary condition on $X$ and Sommerfeld radiation condition at infinity.

\subsection{Some facts of direct scattering}\label{subsec:3.4}

We consider equation \eqref{eq:1.4} assuming for simplicity that
\begin{equation}\label{eq:3.9}
\begin{aligned}
&v \in L^{\infty}(\Omega),\quad v \equiv 0\ \mbox{on}\ \RR^3\setminus \overline {\Omega}\\
&\Omega\ \mbox{is\ an\ open\ bounded\ connected\ domain\ in}\ \RR^3.
\end{aligned}
\end{equation}

The outgoing Green function $R_v^+$ for equation \eqref{eq:1.4} satisfies the integral equation
\begin{equation}
\label{eq:3.10}
R_v^+(x, y, \kappa)=-G^+(x-y,\kappa) +\int_{\Omega}G^+(x-z,\kappa)v(z)R_v^+(z, y, \kappa)dz,
\end{equation}
where $x,y \in \RR^3$, $G^+$ is given by \eqref{eq:3.8b}.

Actually, in addition to \eqref{eq:3.9}, we assume that, for fixed $\kappa>0$,
\begin{equation}
\label{eq:3.11}
\mbox{equation}\ \eqref{eq:3.10}\ \mbox{is\ uniquely\ solvable\ for}\ R_v^+(\cdot , y, \kappa) \in L^2(\Omega).
\end{equation}
In particular, it is known that if $v$ satisfies \eqref{eq:3.9} and is real-valued (or $\Im v \leq 0$), then~\eqref{eq:3.11} is fulfilled automatically;
see, for example, \cite{CK}.

We also consider the scattering wave functions $\psi^+$ for the homogeneous equation~\eqref{eq:1.4} (\ie without $\delta$):
\begin{equation}
\label{eq:3.12}
\psi^+=\psi^+(x,\theta,\kappa)=e^{i \kappa\theta x}+\psi^+_{sc}(x,\theta,\kappa),\quad x \in \RR^3,\ \theta \in \mathbb S^2,
\end{equation}
where $\psi^+_{sc}$ satisfies the radiation condition \eqref{eq:1.2} at fixed $\theta$.

The following formulas hold:
\begin{align}
\label{eq:3.13}
R_v^+(x, y, \kappa)&=R_v^+(y, x, \kappa),\quad x, y\in \RR^3;
\\
\label{eq:3.14}
R_v^+(x, y, \kappa)&=\frac{e^{i\kappa|x|}}{4\pi|x|}\,\psi^+\Bigl(y, -\frac{x}{|x|},\kappa\Bigr)+ O\Bigl(\frac{1}{|x|^2}\Bigr) \quad \mbox{as} \ |x|\ra +\infty\text{ at fixed $y$;}\\
\label{eq:3.15}
\psi^+_{sc}(x,\theta,\kappa)&=\frac{e^{i\kappa|x|}}{|x|}\,A\Bigl(\theta, \frac{x}{|x|},\kappa\Bigr)+ O\Bigl(\frac{1}{|x|^2}\Bigr) \quad \mbox{as} \ |x|\ra +\infty\text{ at fixed $\theta$,}
\end{align}
where $A$ arising in \eqref{eq:3.15} is the scattering amplitude for the homogeneous equation~\eqref{eq:1.4}
and is defined on $\mathbb S^2 \times \mathbb S^2$ at fixed $\kappa$.

In view of \eqref{eq:1.7}, \eqref{eq:1.8}, \eqref{eq:3.12}--\eqref{eq:3.14}, we also have that
\begin{align}
\label{eq:3.16}
R_{v,sc}^+(x, y, \kappa)&=R_{v,sc}^+(y, x, \kappa),\ x, y\in \RR^3;
\\
\label{eq:3.17}
R_{v,sc}^+(x, y, \kappa)&=\frac{e^{i\kappa|x|}}{4\pi|x|}\,\psi_{sc}^+\Bigl(y, -\frac{x}{|x|},\kappa\Bigr)+ O\Bigl(\frac{1}{|x|^2}\Bigr) \ \mbox{as} \ |x|\ra +\infty\text{ at fixed $y$.}
\end{align}
In connection with aforementioned facts concerning $R_v^+$ and $\psi^+$ see, e.g., \cite[Chap.\,IV, \S1]{FM}.

\begin{remark}\label{rem:5}
It is well known that, under assumptions \eqref{eq:3.9}, \eqref{eq:3.11}, the scattering amplitude $A=A(\cdot, \cdot, \kappa)$ is real analytic on $\mathbb S^2 \times \mathbb S^2$.
\end{remark}

\section{Proof of Theorem \ref{th:1}}\label{sec:4}

\subsection{Case $L\subseteq L_{0, \theta}$}\label{subsec:4.1}
First, we give the proof for the case when $L\subseteq L_{0, \theta}$.
In this case it is essentially sufficient to prove that $\Im\psi$ on $L$ uniquely determines $f_j(\theta)$ in~\eqref{eq:3.2} for all $j$.
Such a determination is presented below in this subsection.
The rest follows from the convergence of the series in \eqref{eq:3.2} and analyticity of $\psi$ and $\Im\psi$ on~$L$.

The determination of $f_1$ follows from \eqref{eq:3.5}.
Suppose that $f_1,\dots,f_n$ are determined, then the determination of $f_{n+1}$ is as follows.
Let
\begin{align}
\label{eq:4.1}
\psi_n(x)&=\frac{e^{i\kappa|x|}}{|x|}\sum\limits_{j=1}^{n}\frac{f_j(\theta)}{|x|^{j-1}}, \quad \text{where}\ \theta=\frac{x}{|x|},
\\
\label{eq:4.2}
I_n(x)&=|x|\Im\psi_n,
\\
\label{eq:4.3}
J_n(x)&=|x|^n(I(x)-I_n(x)),
\end{align}
where $x$ is as in \eqref{eq:3.2}, $I(x)$ is defined by \eqref{eq:3.3}.

We have that
\begin{align}
\label{eq:4.4}
2iI(x)&=2iI_n(x)+ \frac{e^{i\kappa|x|}}{|x|^n}{f_{n+1}(\theta)}- \frac{e^{-i\kappa|x|}}{|x|^n}\,\overline{f_{n+1}(\theta)}+O(|x|^{-n-1}),
\\
\label{eq:4.5}
2iJ_n(x)&=e^{i\kappa|x|}{f_{n+1}(\theta)} - e^{-i\kappa|x|}\overline{f_{n+1}(\theta)}+O(|x|^{-1}),
\end{align}
as $|x|\to +\infty$, where $I$ is defined by \eqref{eq:3.3}.

Due to \eqref{eq:4.5} and Remark \ref{rem:4}, we get
\begin{equation}\label{eq:4.6}
\begin{aligned}
&f_{n+1}(\theta)= \frac{1}{\sin(\kappa\tau)}\bigl(-e^{-i\kappa|y|}J_n(x)+e^{-i\kappa|x|}J_n(y)+O(|x|^{-1})\bigr),\\
&x,y\in L_{x_0, \theta},\quad x_0=0,\quad y=x+\tau\theta, \quad \theta\in\mathbb S^2,\quad \tau>0,
\end{aligned}
\end{equation}
assuming that $\sin(\kappa\tau)\ne 0$ for fixed $\tau$
(where the parameter $\tau$ can be always fixed in such a way for fixed $\kappa>0$).

Formulas \eqref{eq:3.3}, \eqref{eq:4.1}--\eqref{eq:4.3} and \eqref{eq:4.6} determine $f_{n+1}$, give the step of induction for finding all $f_j$, and complete the proof of Theorem \ref{th:1} for the case $L\subseteq L_{0, \theta}$.

\subsection{General case}
The general case reduces to the case of Section \ref{subsec:4.1} by the change of variables
\begin{equation}
x'=x-q \label{eq:4.7}\quad
\text{for some fixed $q\in\RR^3$ such that $L\subseteq L_{q, \theta}$.}
\end{equation}
In the new variables $x'\in\mathcal{U}'=\mathcal{U}-q$, we have that:
\begin{equation}
\label{eq:4.8}
\psi= \frac{e^{i\kappa|x'|}}{|x'|}\sum\limits_{j=1}^{\infty}\frac{f_j'(\theta)}{|x'|^{j-1}} \quad \text{for}\ x'\in \RR^3\setminus B_{r'}, \ \theta=\frac{x'}{|x'|},
\end{equation}
for some new $f_j'$,
where $r'$ is such that $\RR^3\setminus B_{r'} \subset \mathcal{U}'$;
\begin{equation}
\label{eq:4.9}
L\subseteq L_{q, \theta}=L_{0, \theta}.
\end{equation}

In addition, the series in \eqref{eq:4.8} converges absolutely and uniformly.

In view of \eqref{eq:4.8}, \eqref{eq:4.9},
we complete the proof of Theorem \ref{th:1} by repeating the proof of Section \ref{subsec:4.1}.

\begin{remark}\label{rem:6}
Our proof of Theorem \ref{th:1} has a holographic prototype in \cite{N5}.
Additional formulas for finding $f_j$ from $\Im\psi$ on $L$
can be obtained proceeding also from the approaches of \cite{N4}, \cite{NS2}, \cite{NSS}.
\end{remark}

\section{Proof of Theorem \ref{th:3}}\label{sec:5}

Under our assumptions on $v$, $\Omega$, $X$, and $\kappa$, we have, in particular, that
\begin{align}
&R_{v,sc}^+(\cdot, y, \kappa)\text{ is real-analytic on}\ X\text{ for fixed}\ y \in X, \label{eq:5.1a} \\
&R_{v,sc}^+(x, \cdot, \kappa)\text{ is real-analytic on}\ X\text{ for fixed}\ x \in X, \label{eq:5.1b}
\end{align}
where $R_{v,sc}^+$ is defined by \eqref{eq:1.7}, \eqref{eq:1.8}.
Here, we use that $R_{v,sc}^+(x, y, \kappa)$ satisfies the homogeneous equation \eqref{eq:1.4} and, therefore, is real-analytic outside of $\supp v$,
and that symmetry \eqref{eq:3.16} holds.

In view of \eqref{eq:1.7}, \eqref{eq:1.8}, \eqref{eq:5.1a}, \eqref{eq:5.1b}, Theorem \ref{th:3} reduces to the case when $\mathcal{D}=X$.
In turn, Theorem \ref{th:3} with $\mathcal{D}=X$ follows from formulas \eqref{eq:1.7}, \eqref{eq:1.8} and from Lemma \ref{lem:1} and Proposition~\ref{prop:2} given below.

\begin{lemma}\label{lem:1}
Under the conditions of Theorem \ref{th:3}, $\Im R_{v,sc}^+(\cdot, y, \kappa)$ on $X$ uniquely determines $R_{v,sc}^+(\cdot, y, \kappa)$ on $X$,
where $y \in \RR^3$, and $\Im R_{v,sc}^+(\cdot, \cdot, \kappa)$ on $X\times X$ uniquely determines $R_{v,sc}^+(\cdot, \cdot, \kappa)$ on $X\times X$.
\end{lemma}

Lemma \ref{lem:1} follows from Theorem \ref{th:2} with $\mathcal U$ such that $\overline {\Omega} \subset \RR^3\setminus \overline {\mathcal{U}}$, $X \subset \mathcal U$, and
the property that $\psi=R_{v,sc}^+(x, y, \kappa)$ is a radiation solution of equation \eqref{eq:1.1} for each $y\in \RR^3$.

\begin{proposition}\label{prop:2}
Under the conditions of Theorem \ref{th:3}, $R_{v,sc}^+(\cdot, \cdot, \kappa)$ on $X\times X$ uniquely determines $v$ (for fixed $\kappa$).
\end{proposition}

Proposition \ref{prop:2} is proved as follows (for example).
First,
\begin{equation}\label{eq:5.2}
\begin{aligned}
&\psi=R_{v,sc}^+(\cdot, x', \kappa)\text{ on}\ X\text{ uniquely determines}\ \psi=R_{v,sc}^+(\cdot, x', \kappa)\text{ on}\ V_X\\
&\mbox{\rm via formula}\ \eqref{eq:3.8a}, \text{ for each fixed }\ x' \in X.
\end{aligned}
\end{equation}

Second,
\begin{equation}\label{eq:5.3}
\begin{aligned}
&R_{v,sc}^+(\cdot, x', \kappa)\text{ on}\ V_X \cup X\text{ uniquely determines}\ \psi_{sc}^+(x',\theta,\kappa)\text{ for}\ \theta \in \Theta_X^+\\
&\mbox{\rm via formula}\ \eqref{eq:3.17}, \text{ for each fixed}\ x' \in X,
\end{aligned}
\end{equation}
where
\begin{equation}
\label{eq:5.4}
\Theta_X^{\pm}=\{\theta \in\mathbb S^2:\ \pm\theta\nu \geq 0 \},
\end{equation}
where $\nu$ is the outward normal to $X$ relative to $V_X$ considered as the interior of $X$.

Third,
\begin{equation}
\begin{aligned}
&\psi=\psi_{sc}^+(\cdot, \theta,\kappa)\text{ on}\ X\text{ uniquely determines}\ \psi=\psi_{sc}^+(\cdot, \theta,\kappa)\text{ on}\ V_X \label{eq:5.5}\\
&\text{via formula \eqref{eq:3.8a}}.
\end{aligned}
\end{equation}

Fourth,\vspace*{-3pt}
\begin{equation}\label{eq:5.6}
\begin{aligned}
&\psi_{sc}^+(\cdot, \theta,\kappa)\text{ on}\ V_X \cup X\text{ uniquely determines}\ A(\theta, \theta', \kappa)\text{ for}\ \theta' \in \Theta_X^-\\
&\text{via formula \eqref{eq:3.15},  for each fixed $\theta \in\mathbb S^2$,}
\end{aligned}
\end{equation}
where $\Theta_X^-$ is defined in \eqref{eq:5.4}.

Therefore, we get that\vspace*{-3pt}
\begin{equation}\label{eq:5.7}
\begin{aligned}
&R_{v,sc}^+(\cdot, \cdot, \kappa)\text{ on}\ X\times X\text{ uniquely determines}\ A(\cdot, \cdot, \kappa)\text{ on}\ \Theta_X^+ \times \Theta_X^-\\
&\text{via \eqref{eq:5.2}, \eqref{eq:5.3}, \eqref{eq:5.5}, \eqref{eq:5.6}.}
\end{aligned}
\end{equation}
In addition,\vspace*{-3pt}
\begin{equation}\label{eq:5.8}
\begin{aligned}
&A(\cdot, \cdot, \kappa)\text{ on}\ \Theta_X^+ \times \Theta_X^-\text{ uniquely determines}\ A(\cdot, \cdot, \kappa)\text{ on}\ \mathbb S^2 \times \mathbb S^2 \\
&\text{by analyticity},
\end{aligned}
\end{equation}
in view of Remark \ref{rem:5}.

Finally, Proposition \ref{prop:2} follows from \eqref{eq:5.7}, \eqref{eq:5.8} and the result of \cite{N1988} (see also \cite{N1994}, \cite{HH}, \cite{E}) that,
under assumptions \eqref{eq:3.9}, \eqref{eq:3.11}, the scattering amplitude $A$ at fixed $\kappa$ uniquely determines $v$.
In this result the assumption that $v$ is real-valued or $\Im v \leq 0$ is not essential and can be replaced by assumption~\eqref{eq:3.11}.

This completes the proofs of Proposition \ref{prop:2} and Theorem \ref{th:3}.

\section{Proof of Theorem \ref{th:4}}\label{sec:6}
Under the assumptions of Theorem \ref{th:4}, the functions $\psi$ and $\Im \psi$ are real-analytic in $\mathcal{U}$, in general, and on $Y$, in particular.
Therefore, $\Im \psi$ on $\mathcal{D}$ uniquely determines $\Im \psi$ on $Y$ by analytic continuation, taking also into account that $Y$ is real-analytic and connected.
In turn, $\Im \psi$ on $Y$ uniquely determines $\Im \psi$ in $\mathcal{B}$ by solving the Dirichlet problem for the Helmholtz equation.
In turn, $\Im \psi$ in $\mathcal{B}$ uniquely determines $\Im \psi$ in $\mathcal{U}$ by analytic continuation.
Finally, $\Im \psi$ in $\mathcal{U}$ uniquely determines $\psi$ in $\mathcal{U}$, in~view of Corollary \ref{cor:1} in Section \ref{sec:2}.

This completes the proof of Theorem \ref{th:4}.

\section{Proof of Theorem \ref{th:5}}\label{sec:7}
Recall that $\Im R_v^+(\cdot, \cdot, \kappa)$ on $\mathcal{D} \times \mathcal{D}$ uniquely determines $\Im R_{v,sc}^+(\cdot, \cdot, \kappa)$ on $\mathcal{D} \times \mathcal{D}$,
in view of \eqref{eq:1.7}, \eqref{eq:1.8}.
Under the assumptions of Theorem \ref{th:5}, we have, in particular, that
\begin{equation}
\label{eq:7.1}
\psi=R_{v,sc}^+(\cdot, y, \kappa)\text{ is a radiation solution of equation \eqref{eq:1.1}},\ y\in \RR^3.
\end{equation}
Therefore, due to Theorem \ref{th:4}, $\Im R_{v,sc}^+(\cdot, y, \kappa)$ on $\mathcal{D}$ uniquely determines $R_{v,sc}^+(\cdot, y, \kappa)$ in $\mathcal{U}$.

Similarly, $\Im R_{v,sc}^+(x, \cdot, \kappa)$ on $\mathcal{D}$ uniquely determines $R_{v,sc}^+(x, \cdot, \kappa)$ in $\mathcal{U}$,
due to \eqref{eq:3.16}, where $x\!\in\! \mathbb R^3$.
Therefore, $\Im R_{v,sc}^+(\cdot, \cdot, \kappa)$ on $\mathcal{D} \times \mathcal{D}$ uniquely determines~$R_{v,sc}^+(\cdot, \cdot, \kappa)$ on $\mathcal{U} \times \mathcal{U}$.
Finally, $R_{v,sc}^+(\cdot, \cdot, \kappa)$ on $\mathcal{U} \times \mathcal{U}$ uniquely determines $v$ by different ways.

For example: $R_{v,sc}^+(\cdot, \cdot, \kappa)$ on $\mathcal{U} \times \mathcal{U}$ uniquely determines
$A(\cdot, \cdot, \kappa)$ on $\mathbb S^2 \times \mathbb S^2$, in~view of~\eqref{eq:3.15}, \eqref{eq:3.17};
$A$ at fixed $\kappa$ uniquely determines $v$ as recalled at the end of proof of Theorem \ref{th:3}.

This completes the proof of Theorem \ref{th:5}.

\begin{remark}\label{rem:7}
Formulas relating $R_v^+$ on $\mathbb S^2_r \times \mathbb S^2_r$ and $A$ on $\mathbb S^2 \times \mathbb S^2$ at fixed $\kappa$,
where $v(x) = 0$ for $|x|>r$, $\mathbb S^2_r$ is defined by \eqref{eq:2.1}, were given for the first time in \cite{Ber}.
\end{remark}

\begin{remark}\label{rem:8}
The case is also of interest when the boundary $Y$ in \eqref{eq:2.2} is not connected but consists of two disjoint connected components
$Y_1$ and $Y_2$, where $Y_1=\partial \mathcal{B}_1$, $Y_2=\partial \mathcal{B}_2$ and $\mathcal{B}_1$, $\mathcal{B}_2$ are open bounded domains such that $\RR^3 \setminus \mathcal{U} \subset \mathcal{B}_1 \subset \mathcal{B}_2$.
In this case Theorem \ref{th:4} is valid with $\mathcal{D}$ replaced by $\mathcal{D}_1 \cup \mathcal{D}_2$, whereas Theorem \ref{th:5} is valid with $\mathcal{D} \times \mathcal{D}$ replaced by $(\mathcal{D}_1 \cup \mathcal{D}_2) \times \mathcal{D}_1$
(as well as with $\mathcal{D} \times \mathcal{D}$ replaced by $(\mathcal{D}_1 \cup \mathcal{D}_2) \times \mathcal{D}_2$, where $\mathcal{D}_1$, $\mathcal{D}_2$ are arbitrary non-empty open domains of $Y_1$ and $Y_2$, respectively.
It~is of interest to compare the later result with Theorem \ref{th:4}.3 in the recent thesis \cite{Mu}, which gives uniqueness for monochromatic passive imaging from cross correlations on $(Y_1 \cup Y_2) \times (Y_1 \cup Y_2)$.
These comparisons may lead to further important results.
\end{remark}

\backmatter
\bibliographystyle{jepplain+eid}
\bibliography{novikov}
\end{document}
